\input pdfToolbox
\input pdfmsym
\pdfmsymsetscalefactor{10}
\input ../preamble
\def\enumindent{0cm}

\def\printmcount{\the\counter{section}.\the\counter{math counter}}
\setcounter{section}{9}

\centerline{\setfontandscale{bf}{20pt}Commutative Algebra \& Algebraic Geometry}
\medskip
\centerline{\bf Exercise \the\counter{section}}
\smallskip
\centerline{\it Ari Feiglin}

\def\h{{\sf h}} \def\dh{{\sf dh}}
\def\rad{{\rm rad}}

\bprob

    \benum
        \item Prove that the homogenization and dehomogenization of polynomials are inverses, except for losing the highest $x_0$ factor.
        \item Consider $I=(f,g)\leq k[x,y,z]$ where $f=y-x^2,g=z-x^3$.
        Show that the homogenization $I_h\leq k[x,y,z,w]$ does not equal $(f_h,g_h)$ by showing that $xy-zw\in I_h-(f_h,g_h)$.
        \item Let $f\in k[x_0,\dots,x_n]$ be an irreducible homogeneous polynomial.
        Show that the degree of $Z_{\bP^n}(f)$ is independent of $f$.
    \eenum

\eprob

\benum
    \item For a multi-index $\alpha=(\alpha_1,\dots,\alpha_n)$ and indeterminants $x=(x_1,\dots,x_n)$ define
    $$ x^\alpha = \prod_{i=1}^n x_i^{\alpha_i} . $$
    Let $x/x_0=(x_1/x_0,\dots,x_n/x_0)$, then for $f\in k[x_1/x_0,\dots,x_n/x_0]$ we can write it as
    $$ f = \sum_\alpha a_\alpha (x/x_0)^\alpha, $$
    where the sum ranges over a finite set of multi-indices and $a_\alpha\in k$.
    We can then order the sum by degree of multi-index (where $\abs\alpha=\sum\alpha_i$)
    $$ f = \sum_{n=0}^d\sum_{\abs\alpha=n}a_\alpha(x/x_0)^\alpha = \sum_{n=0}^d\sum_{\abs\alpha=n}a_\alpha\frac{x^\alpha}{x_0^n} . $$
    Thus $f$'s homogenization is
    $$ f_\h = \sum_{n=0}^d\sum_{\abs\alpha=n}a_\alpha x_0^{d-n}x^\alpha . $$
    The dehomogenization of this is
    $$ (f_\h)_\dh = f_\h(1,x/x_0) = \sum_{n=0}^d\sum_{\abs\alpha=n}a_\alpha(x/x_0)^\alpha = f , $$
    as desired.

    Now for $g\in k[x_0,\dots,x_n]$ homogeneous of degree $d$ we write
    $$ g = \sum_\beta b_\beta\hat x^\beta $$
    for $b_\beta\in k$, $\hat x=(x_0,\dots,x_n)$, and all the degrees of the multi-indices in the sum are equal to $d$.
    For multi-index $\beta$, we can write it as $\beta=(n,\alpha)$ for multi-index $\alpha$ over $(x_1,\dots,x_n)$.
    Since $\abs\beta=n+\abs\alpha=d$, we have that $n=d-\abs\alpha\leq d$.
    So we can then rewrite $g$ as (writing $b_\alpha=b_{(d-n,\alpha)}$):
    $$ g = \sum_{n=0}^d\sum_{\abs\alpha=n} b_{\alpha}x_0^{d-n}x^\alpha . $$
    Its dehomogenization is
    $$ g_\dh = g(1,x/x_0) = \sum_{n=0}^d\sum_{\abs\alpha=n} b_\alpha\frac{x^\alpha}{x_0^n} . $$
    This is a polynomial whose degree is the maximal $N$ such that there is a multi-index where $\abs\alpha=N$ and $b_{\alpha}\neq0$.
    Then
    $$ (g_\dh)_\h = \sum_{n=0}^d\sum_\alpha b_{\alpha}x_0^{N-n}x^\alpha . $$
    So we see that
    $$ g = x_0^{d-N}(g_\dh)_\h , $$
    and $x_0$ does not divide $(g_\dh)_\h$, since for $\alpha$ of degree $N$ with $b_\alpha\neq0$, in $(g_\dh)_\h$ its component is
    $b_\alpha x^\alpha$ which is not divided by $x_0$; so we have removed the highest factor of $x_0$ from $g$.

    \item Firstly, we see
    $$ (xf-g)_\h = (xy-x^3 - z + x^3)_\h = (xy - z)_\h = xy - zw , $$
    and since $xf-g\in I$ we see that $xy-zw\in I_\h$.

    We see that
    $$ f_\h = wy - x^2,\qquad g_\h = w^2z - x^3 . $$
    And $xy-zw$ cannot be in the ideal generated by these.
    Indeed, multiplying any monomial by $f_\h$ or $g_\h$ cannot yield an $xy$ term: $X^\alpha wy$ and $X^\alpha x^2$ both are not $xy$ for example.
    Thus multiplying any polynomial by $f_\h$ or $g_\h$ yields a sum of monomials, none of which are $xy$, and so the sum itself cannot yield an $xy$ term.
    And so the sum of products of polynomials with $f_\h,g_\h$ cannot yield an $xy$ term, and thus cannot be equal to $xy-zw$.

    \item Let $f,g\in k[x_0,\dots,x_n]$ be irreducible homogeneous polynomials with $Z_{\bP^n}(f)=Z_{\bP^n}(g)$.
    We want to show that they have the same degree.
    In question 2 we will show that if $Z_{\bP^n}(f)$ is nonempty then $\I(Z_{\bP^n}(f))=\sqrt{(f)}$.
    Since $f$ is irreducible, this is equal to $(f)$.
    So in the case that the zerosets are nonempty we have $(f)=(g)$ so $f=gp$ and $g=fq$ for polynomials $p,q$.
    But then $\deg(f)=\deg(g)+\deg(q)$ and $\deg(g)=\deg(f)+\deg(q)$, meaning $\deg(f)\leq\deg(g)\leq\deg(f)$ so we have equality.
\eenum

\bprob

    Let $n\in\bN$.
    \benum
        \item The {\emphcolor irrelevant ideal} of $k[x_0,\dots,x_n]$ is
        $$ k[x_0,\dots,x_n]_+=\bigoplus_{d>0}k[x_0,\dots,x_n]_d . $$
        Let $I\leq k[x_0,\dots,x_n]$ be a homogeneous ideal.
        Show that the following are equivalent:
        $$ (1)\ Z_{\bP^n}(I)=\emptyset,\qquad (2)\ k[x_0,\dots,x_n]_+ \leq \sqrt I,\qquad
        (3)\ k[x_0,\dots,x_n]_d \subseteq I\hbox{ for some $d$} . $$
        Such an ideal is called an {\emphcolor irrelevant ideal} (as opposed to {\it the\/} irrelevant ideal).
        \item Let $J\leq k[x_0,\dots,x_n]$ be a homogeneous ideal which is not irrelevant.
        Show that $\I(\bZ_{\bP^n}(J))=\sqrt J$.

        Conversely show that for $Y\subseteq\bP^n$ closed, $Z_{\bP^n}(\I(Y))=\cl Y$.
        \item Let $\cll(\bP^n)$ be the closed subsets of $\bP^n$ and $\rad_h$ be the radical homogeneous non-irrelevant ideals of
        $k[x_0,\dots,x_n]$.
        We have maps
        $$ \I\colon\cll(\bP^n) \tof \rad_h\colon Z . $$
        Show the following.
        \benum
            \item $\I,Z$ form an antiequivalence of posets.
            \item They send irreducible closed subsets to prime ideals and vice versa.
            \item The maps interchange points $[a_0:\cdots:a_n]$ and ideals $\gen{a_ix_j-a_jx_i}_{i,j}$.
        \eenum

        \item Show that $\bP^n$ is irreducible.

        \item Show that $\bP^n$ is a Noetherian topological space of dimension $n$.
    \eenum

\eprob

\benum
    \item Since the projection $\pi\colon\bA^{n+1}-0\to\bP^n$ is surjective, $Z_{\bP^n}(I)$ is empty iff its preimage is.
    Thus
    $$ Z_{\bP^n}(I) = \emptyset \iff \emptyset = \pi^{-1}Z_{\bP^n}(I) = Z_{\bA^{n+1}-0}(I) \iff Z_{\bA^{n+1}}(I) \subseteq \set0 . $$
    We know that $\set 0 = Z_{\bA^{n+1}}(x_0,\dots,x_n)$, so this is iff
    $$ \sqrt{(x_0,\dots,x_n)} \leq \sqrt I , $$
    and $(x_0,\dots,x_n)$ is maximal so it is radical.
    Furthermore, $k[x_0,\dots,x_n]_+=(x_0,\dots,x_n)$ because the irrelevant ideal clearly contains this maximal ideal as $x_i$ are homogeneous
    of degree $1$, and the irrelevant ideal is proper so there is equality by maximality.

    So we have shown $Z_{\bP^n}(I)=\emptyset$ iff $k[x_0,\dots,x_n]_+\leq\sqrt I$, which is the equivalence of (1) and (2).

    If $k[x_0,\dots,x_n]_d\leq I$ for any $d$ then $Z_{\bP^n}(I)$ is empty.
    This is because it implies $x_0^d,\dots,x_n^d\in I$ and their only common zero is $0$.
    So $Z_{\bA^{n+1}}(I)=\set0$ and thus $Z_{\bP^n}(I)=\emptyset$.
    So (3) implies (1).

    Conversely suppose (1), so we get (2) by equivalence.
    So $x_0,\dots,x_n\in\sqrt I$ and thus there are $k_0,\dots,k_n$ such that $x_0^{k_0},\dots,x_n^{k_n}\in I$.
    Let $d=k_0+\cdots+k_n$, then $k[x_0,\dots,x_n]_d\subseteq I$.
    Indeed if $\alpha$ is any multi-index of degree $d$, then there is some $\alpha_i\geq k_i$ so $x_i^{k_i}\in I$ occurs in the product
    for $x^\alpha$, so $x^\alpha\in I$.
    So then any linear combination of monomials of degree $d$, i.e.\ an element of $k[x_0,\dots,x_n]_d$, is in $I$ as desired.

    \item Since $J$ is homogeneous, $J=\bigoplus_d k[\bar x]_d\cap J$, and since it is not irrelevant it cannot contain any scalars (as then
    it would have empty zero set), so this sum can be taken over $d>0$ showing that $J$ is contained in the irrelevant ideal.
    Thus $J\leq(x_0,\dots,x_n)$. 

    Now we claim
    $$ \I(Z_{\bP^n}(J)) = \I_\bA(Z_{\bA^{n+1}}(J)) = \sqrt J . $$
    The left-hand side is generated by homogeneous $f$ such that for every $P\in Z_{\bP^n}(J)$, $f(P)=0$.
    Since $Z_{\bP^n}(J)$ is nonempty these $f$ are not scalars (the condition cannot be vaccuous).
    This means that for every lift of $P$ to $\tilde P$, $f(\tilde P)=0$.
    For $\tilde P\in Z_{\bA^{n+1}}(J)$, if $\tilde P\neq0$ then $\pi(\tilde P)\in Z_{\bP^n}(J)$ as we showed in class.
    Thus $f(\tilde P)=0$ so $f\in\I_\bA(Z_\bA(J))$.
    And for $\tilde P=0$, since $f$ is homogeneous and not a scalar, $f(0)=0$.
    Thus $f(\tilde P)=0$ in all cases, so $f\in\I_\bA(Z_\bA(J))$, showing left-to-right inclusion.

    Now $Z_{\bA^{n+1}}(J)$ is invariant under $k^\times$'s action, since $J$ is generated by homogeneous polynomials and their roots
    are invariant under $k^\times$.
    As we showed in class, this means $\I_\bA(Z_\bA(J))$ is homogeneous and thus generated by its homogeneous elements.
    For any homogeneous $f$ in this ideal, it is also in $\I(Z_\bP(J))$: take any $P\in Z_\bP(J)$ and lift it to $\tilde P$.
    Then $\tilde P\in\pi^{-1}Z_\bP(J)=Z_\bA(J)-0$ and so $f(\tilde P)=0$ so $f(P)=0$.
    So all of $\I_\bA(Z_\bA(J))$'s homogeneous elements are in $\I(Z_\bP(J))$ and as it is generated by those, there is inclusion and thus
    equality.

    For the converse, closed sets of $\bP^n$ are of the form $Z_\bP^n(J)$ for homogeneous ideals $J$.
    Let $Y\subseteq\bP^n$ be nonempty, and suppose $Y\subseteq Z_\bP(J)$; since $Y$ is nonempty $J$ is not irrelevant.
    Then since $\I$ is order-reversing, we see that from what we just proved
    $$ \I(Y) \supseteq \I(Z_\bP(J)) = \sqrt J \supseteq J . $$
    And so
    $$ Z_\bP(\I(Y)) \subseteq Z_\bP(J) . $$
    And $Y\subseteq Z_\bP(\I(Y))$ clearly, so it is the smallest closed set containing $Y$; its closure.

    \item \benum
        \item This is clear, as by the previous points $Z$ maps non-irrelevant ideals to non-empty closed subsets, and thus
        $$ \I(Z(J)) = \sqrt J = J $$
        for non-irrelevant radical $J$.
        And $Z(\I(Y))=\cl Y$ so for $Y$ closed this is just $Y$.

        \item For $\p$ prime, suppose $Z(\p)$ is the union of two closed subsets, $Z(\p)=Z(I)\cup Z(J)=Z(IJ)$.
        Then
        $$ \p = \I(Z(\p)) = \I(Z(IJ)) = \sqrt{IJ} , $$
        where the final equality is because $Z(IJ)$ is nonempty (as it is equal to $Z(\p)$).
        In particular we have $IJ\leq\p$, and since $\p$ is prime one of them is contained in $\p$.
        Say $I\leq\p$ so $Z(I)\supseteq Z(\p)$ and thus we have equality.
        So $Z(\p)$ is irreducible.

        Now suppose $Z(J)$ is irreducible for $J$ radical; we claim that $J$ is prime.
        Suppose $I_0I_1\leq J$ for ideals $I_0,I_1$, then $Z(J)\subseteq Z(I_0I_1)=Z(I_0)\cup Z(I_1)$.
        Since $Z(J)$ is irreducible this means that $Z(J)\subseteq Z(I_0)$ or $Z(I_1)$, say $Z(I_0)$.
        This implies, as $Z(I_0)$ is then non-empty,
        $$ I_0 \subseteq \sqrt{I_0} = \I(Z(I_0)) \leq \I(Z(J)) = \sqrt J = J . $$
        Thus $J$ is prime, as desired.
        Since $\I(Z(J))=J$, this shows that $\I$ maps irreducible closed sets to prime ideals.

        \item Choose $i$ such that $a_i\neq0$.
        Then $x$ is a zero for $a_ix_j-a_jx_i$ for all $j$ means that $x_j=(a_j/a_i)x_i$.
        Thus
        $$ [x_0:\cdots:x_n] = \bracks{\frac{a_0}{a_i}:\cdots:\frac{a_n}{a_0}} = [a_0:\cdots:a_n] , $$
        so that
        $$ Z\gen{a_ix_j-a_jx_i} = \set{[a_0:\cdots:a_n]} , $$
        since clearly $(a_0,\dots,a_n)$ is a root of each $a_ix_j-a_jx_i$.
        We claim that $I=\gen{a_ix_j-a_jx_i}$ is prime; this is sufficient because then as $\I,Z$ are inverses on non-irrelevant radical
        ideals this means $\I[a_0:\cdots:a_n]=I$.

        Take $a_i\neq0$ and define
        $$ \phi\colon k[x_0,\dots,x_n] \to k[t],\qquad x_j \mapsto \frac{a_j}{a_i}x_i . $$
        This is surjective as $x_i$ is mapped to $t$.
        Clearly $I$ is contained in the kernel of $\phi$, we claim there is equality.
        Let $f\in\ker\phi$, then modulo $I$ all $x_j$ are scalar multiples of $x_i$, so $f$ is equivalent to $g(x_i)$ modulo $I$ for some
        single-variable $g$.
        Since $I$ is contained in the kernel, this means that $\phi(f)=\phi(g(x_i))=0$, but $\phi(g(x_i))=g(t)$ and this is zero iff
        $g=0$.
        So $f$ is equal to zero modulo $I$, meaning $f\in I$.

        Thus $\ker\phi=I$ and so $\phi$ induces an isomorphism $k[x_0,\dots,x_n]/I\cong k[t]$ which is a domain so $I$ is prime, as desired.
    \eenum

    \item $\bP^n$ is closed and $\I(\bP^n)=0$ since the only polynomial which zeroes all of $\bP^n$ is the zero polynomial.
    The zero ideal is prime, and so by the previous point this means $\bP^n$ is irreducible.

    \item $\bP^n$ is an algebraic variety and those are always Noetherian.
    This is because it has a finite open affine cover, and affine varieties are Noetherian.
    A finite union of open Noetherian subsets is Noetherian: let $X=\bigcup_iU_i$ for $U_i$ Noetherian, then for every $Y\subseteq X$ we have $Y=\bigcup_iU_i\cap Y$.
    Since $U_i$ is Noetherian, $U_i\cap Y$ is quasicompact, and thus $Y$ is quasicompact as a finite union of quasicompact sets.
    So all of $X$'s subsets are quasicompact, meaning $X$ is Noetherian.

    $\bP^n$ has an open affine cover by affine varieties $U_i\cong\bA^n$.
    We showed that for an open cover of a topological space $X=\bigcup_iU_i$, $\dim X=\sup_i\dim U_i$.
    In this case
    $$ \dim\bP^n = \sup_i\dim U_i = \sup_i\dim\bA^n = \sup_in = n , $$
    as desired.
\eenum

\bprob

    Let $X$ be an affine variety and fix $n\in\bN$.
    \benum
        \item Let $f_0,\dots,f_n\in k[X]$.
        Define $X'=X-Z(f_0,\dots,f_n)$, and define $f\colon X'\to\bP^n$ by $f(P)=[f_0(P):\cdots:f(P)]$.
        Prove that this is a well-defined morphism of varieties.

        \item Let $m\in\bN$, $X=\bA^{m+1}$, and let $f_0,\dots,f_n\in k[X]$ be homogeneous of the same degree such that $Z(f_0,\dots,f_n)=\set0$.
        Thus by point (1) we get a morphism $f\colon X-\set0\to\bP^n$.
        Show that $f$ induces a well-defined morphism $\bar f\colon\bP^m\to\bP^n$.
    \eenum

\eprob

\benum
    \item Note that $f$ decomposes as
    $$ f = X' \xtos{f'} \bA^{n+1} - 0 \xtos\pi \bP^n $$
    where $f'(P)=(f_0(P),\dots,f_n(P))$ and $\pi$ is the projection morphism.
    $f'$ is a morphism of varieties since each of its components are regular.
    Explicitly $f'$ is a morphism into $\bA^{n+1}$ because of this, and it maps into $\bA^{n+1}-0$ so by the universal property of the induced
    structure on $\bA^{n+1}-0$ it is a morphism to $\bA^{n+1}-0$.
    Thus $f$ is a composition of morphisms and thus a morphism itself.

    \item Consider the following commutative diagram

    \commsq{\bA^{m+1}-0&\bA^{n+1}-0\cr\bP^m&\bP^n}
        {f'}{\bar f}{\pi}{\pi}

    Where $f'$ is the map from the previous point and $\bar f$ is defined as
    $$ \bar f(P) = [f_0(P):\cdots:f_n(P)] . $$
    This is well-defined because for any two lifts of $P$, $\tilde P_1=\lambda\tilde P_0$ we have that
    $$ [f_0(\tilde P_1):\cdots:f_n(\tilde P_1)] = [\lambda^df_0(\tilde P_0):\cdots:\lambda^df_n(\tilde P_0)] = [f_0(\tilde P_0):\cdots:f_n(\tilde P_0)] , $$
    where $d$ is the common degree of the $f_i$s.
    The diagram clearly commutes: either path gives $P\mapsto[f_0(P):\cdots:f_n(P)]$ for $P\in\bA^{m+1}-0$.

    Thus $\bar f\circ\pi=\pi\circ f'=f$, and since $\pi$ is the quotient map, by its universal property it means that $\bar f$ is a morphism of varieties.
\eenum

\bprob

    Let $f=yz^2-x^3+3xz^2-z^3$ and $g=yz-x^2+2xz$ in $\bC[x,y,z]$.
    Denote $C=Z_\bP(f),D=Z_\bP(g)$.
    Prove Bezout's theorem by hand in this case.

\eprob

We first find $C\cap D=Z_\bP(f,g)$.
Note
$$ (f,g) = (f-gz,g) = (-(x-z)^2(x+z),g) = ((x-z)^2(x+z), g) . $$
So any (affine) zero has to have $x=\pm z$.
If $x=z$ then we have
$$ 0 = g(z,y,z) = yz + z^2 \implies z = 0 \hbox{ or } y = -z . $$
If $x=-z$ then
$$ 0 = g(-z,y,z) = yz - 3z^2 \implies z = 0 \hbox{ or } y = 3z . $$
So the affine zeroes are
$$ Z_\bA(f,g) = \set{(0,y,0),(z,-z,z),(-z,3z,z)} , $$
and thus projecting onto $\bP^2$ gives
$$ Z_\bP(f,g) = \set{[0:1:0],\ [1:-1:1],\ [-1:3:1]} . $$
So we need to find the multiplicty at each point.

For $P=[0:1:0]$, we take the local affine neighborhood $U_1$, in which $P$ corresponds to the point $(0,0)$, and $f,g$ correspond to their dehomogenizations:
$$ f_\dh = z^2 - x^3 + 3xz^2 - z^3,\qquad g_\dh = z - x^2 + 2xz . $$
And the stalk at $P$ is $k[x,z]_{(x,z)}$.
When quotienting by $g_\dh$, we see that
$$ z = \frac{x^2}{1+2x} , $$
and thus
$$ f_\dh = x^3\parens{\frac x{(2x+1)^2} - 1 + \frac{3x^2}{(2x+1)^2} - \frac{x^3}{(2x+1)^3}} . $$
The term in the parentheses evaluates to $-1$ at $x=0$, which means it is a unit as this is nonzero at $P$.
So
$$ \frac{\O_{\bP^2,P}}{(f,g)} = \frac{k[x]_{(x)}}{(x^3)} . $$
This has a basis $1,x,x^2$ so the multiplicity at $[0:1:0]$ is $3$.

For $P=[1:-1:1]$ we take the local affine neighborhood $U_2$, in which $P$ corresponds to $(1,-1)$, and $f,g$ correspond to their dehomogenizations:
$$ f_\dh = y - x^3 + 3x - 1,\qquad g_\dh = y - x^2 + 2x . $$
Now we can write these in terms of $x=u+1$ and $y=v-1$, so that $P$ then corresponds to $0$ and
$$ f_\dh = v - u^3 - 3u^2,\qquad g_\dh = v - u^2 . $$
Then
$$ (f_\dh,g_\dh) = (u^3+2u^2,v-u^2) , $$
and $u^3+2u^2=u^2(u+2)$ where $u+2$ is a unit because at $P$ it is nonzero, so
$$ (f_\dh,g_\dh) = (u^2,v-u^2) = (u^2,v) . $$
Thus we have
$$ \frac{\O_{\bP^2,P}}{(f,g)} = \frac{k[u,v]_{(u,v)}}{(u^2,v)} , $$
which has basis $1,u$; the multiplicity at $[1:-1:1]$ is then $2$.

For $P=[-1:3:1]$ we take local affine neighborhood $U_2$, in which $P$ corresponds to $(-1,3)$ and $f,g$ correspond to the same dehomogenizations as previously.
Again we shift our coordinates, $x=u-1$ and $y=v+3$ so
$$ f_\dh = v - u^3 + 3u^2,\qquad g_\dh = v - u^2 + 4u . $$
And so
$$ (f_\dh,g_\dh) = (v - u^2 + 4u, u^3 - 4u^2 + 4u) = (v-u^2+4u,u(u^2-4u+4)) = (v-u^2+4u,u) = (v,u) , $$
where the second-to-last equality is because $u^2-4u+4$ is a unit in the stalk.
Thus $(f_\dh,g_\dh)$ is maximal, they intersect transversally at $P$ and so the multiplicity is $1$.

So the total sum of multiplicities is $3+2+1=6$.
And $\deg f=3,\deg g=2$ so we have demonstrated Bezout's theorem.

\bprob

    Fix $n>0$.
    \benum
        \item Let $h\in k[x_0,\dots,x_n]$ be a linear form, i.e.\ a homogeneous polynomial of degree $1$.
        Show that the hyperplane $H=Z_\bP(h)$ is isomorphic to $\bP^{n-1}$.

        \item Fix $P=[a_0:\cdots:a_n]\in\bP^n$ not in $H$.
        Define $\phi\colon\bP^n-P\to H$ as follows.
        For $Q\in\bP^n-P$, connect it to $P$ by a line, then let $\phi(Q)$ be the intersection of that line with $H$.
        Show that
        $$ \phi([b_0:\cdots:b_n]) = \bracks{b_i - \frac{h(b_0,\dots,b_n)}{h(a_0,\dots,a_n)}a_i}_i . $$
    \eenum
    I unfortunately do not have time to draw, but if I did I would draw the empty box of hot chocolate packets lying on my desk\dots

\eprob

\benum
    \item Let us write
    $$ h = \sum_i\beta_ix_i,\qquad \beta_i\in k . $$
    There is some $i$ for which $\beta_i\neq0$, potentially by changing variables we can assume $i=0$.
    Let us write the coordinates of $\bP^{n-1}$ by $[a_1:\cdots:a_n]$.
    Now consider $\Phi\colon\bA^n-0\to\bP^n$ given by question 3 and the polynomials $-\frac1{\beta_0}\sum_{i>0}\beta_ix_i$ and $x_1,\dots,x_n$; we write it as $\Phi=\pi\circ\Phi'$ as
    in question 3.
    These are homogeneous of the same degree so by question 3 part 2, this induces a morphism $\bar\Phi\colon\bP^{n-1}\to\bP^n$:
    $$ \bar\Phi[a_1:\cdots:a_n] = \bracks{-\frac1{\beta_0}\sum_{i>0}\beta_ia_i:a_1:\cdots:a_n} . $$
    Its image is equal to $H$, and it is injective by considering the last $n$ coordinates in the image of $\bar\Phi$.
    So it restricts to a bijective morphism $\bar\Phi\colon\bP^{n-1}\to H$.

    Now, $\Phi'$'s image is $Z_\bA(h)-0$ and it is injective too.
    Its inverse is given by $(\Phi')^{-1}(a_0,\dots,a_n)=(a_1,\dots,a_n)$ which is a morphism of varieties as each component is regular.
    Taking inverses, we have the commutative diagram

    \commsq{\bA^{n-1}-0&Z_\bA(h)-0\cr\bP^n&Z_\bP(h)}
        {(\Phi')^{-1}}{\bar\Phi^{-1}}{\pi}{\pi}

    and since the $\pi$s are quotients, this means $\bar\Phi^{-1}$ is a morphism as well.
    Thus $\bar\Phi$ is an isomorphism as desired.

    \item The line connecting $[b_0:\cdots:b_n]$ and $[a_0:\cdots:a_n]$ is given by the points of the form $[sb_i+ta_i]_i$.
    This is in $H$ iff
    $$ 0 = h[sb_i+ta_i]_i = h(sb_i+ta_i)_i = sh(b) + th(a) , $$
    so we have that $t=-sh(b)/h(a)$ so the resulting point is
    $$ \bracks{sb_i - s\frac{h(b)}{h(a)}}_i = \bracks{b_i - \frac{h(b)}{h(a)}}_i , $$
    as desired.

    The coordinates of $\phi$ are all polynomial in $b$, so it is a morphism of varieties.
\eenum

\bye

